\section{Related Work}
\subsection{Interactive World Models}

Recent advances in video generation models \cite{sora2,longlive,hacohen2026ltx} have enabled the development of interactive world generation models \cite{genie,cosmos,wan2025wan}.
In the realm of video generation,
text-to-video \cite{chen2024videocrafter2,li2024t2v,yume} and image-to-video generation \cite{xing2024dynamicrafter,xu2024easyanimate,shi2024motion,wan2025wan} have achieved remarkable progress in generation quality and temporal consistency.
As for interactive video generation,
works \cite{yume1_5,genie3,team2026advancing} enable interaction by switching prompts during the generation process,
while works \cite{genie3,hyworld2025,gao2025longvie, matrixgame1,matrixgame2,hunyuangame,yan,voyager} introduce actions via keyboard control \cite{matrixgame2,genie3} and camera poses \cite{camctrl1, camctrl2} on top of image-to-video generation to control the generated video.
Despite producing promising results, these methods are limited by a restricted action space and tightly couple action control with pixel-level video changes.
In contrast, we focus on state-aware video generation via action control, which features a rich action space (over 450 actions) and aims to use states as an intermediate representation to convey the effects of actions on pixel-level video generation.

Some recent works \cite{wang2026mechanistic, yue2025simulating, garrido2026learning, lillemarkflow, 3d_as_code} attempt to introduce latent state representations into video generation models to better capture environment dynamics.
However, these approaches typically represent the world state as an implicit latent variable learned from visual observations.
By contrast, we focus on explicit, semantically meaningful states and introduce WildWorld, a large-scale dataset with state annotations for learning and analyzing state dynamics.

\subsection{Video Generation Dataset}

Recent progress in video generation has been driven by several large-scale datasets, such as OpenVid-1M \cite{nan2024openvid}, MiraData \cite{ju2024miradata}, Open-Sora \cite{lin2024open}, and SpatialVID \cite{spatialvid}, which provide large collections of internet videos for training generative models. 
More recent works have begun to explore datasets for world modeling or interactive video generation, including OmniWorld \cite{zhou2025omniworld}, Sekai \cite{sekai}, GF-Minecraft \cite{yu2025gamefactory}, PLAICraft \cite{he2025plaicraft}, and GameGen-X \cite{che2024gamegen}.
While these datasets introduce gameplay videos or action signals to capture environment dynamics, they still primarily rely on visual observations and lack explicit, semantically meaningful state representations.
MIND \cite{ye2026mind} proposes a benchmark for evaluating memory consistency and action control in world models.
Compared to the above works, WildWorld provides explicit state annotations, including character skeletons, world states, camera poses, and depth, enabling models to learn structured state dynamics and supporting direct evaluation of state alignment and action following.